% Options for packages loaded elsewhere
\PassOptionsToPackage{unicode}{hyperref}
\PassOptionsToPackage{hyphens}{url}
\documentclass[
]{article}
\usepackage{xcolor}
\usepackage[margin=1in]{geometry}
\usepackage{amsmath,amssymb}
\setcounter{secnumdepth}{-\maxdimen} % remove section numbering
\usepackage{iftex}
\ifPDFTeX
  \usepackage[T1]{fontenc}
  \usepackage[utf8]{inputenc}
  \usepackage{textcomp} % provide euro and other symbols
\else % if luatex or xetex
  \usepackage{unicode-math} % this also loads fontspec
  \defaultfontfeatures{Scale=MatchLowercase}
  \defaultfontfeatures[\rmfamily]{Ligatures=TeX,Scale=1}
\fi
\usepackage{lmodern}
\ifPDFTeX\else
  % xetex/luatex font selection
\fi
% Use upquote if available, for straight quotes in verbatim environments
\IfFileExists{upquote.sty}{\usepackage{upquote}}{}
\IfFileExists{microtype.sty}{% use microtype if available
  \usepackage[]{microtype}
  \UseMicrotypeSet[protrusion]{basicmath} % disable protrusion for tt fonts
}{}
\makeatletter
\@ifundefined{KOMAClassName}{% if non-KOMA class
  \IfFileExists{parskip.sty}{%
    \usepackage{parskip}
  }{% else
    \setlength{\parindent}{0pt}
    \setlength{\parskip}{6pt plus 2pt minus 1pt}}
}{% if KOMA class
  \KOMAoptions{parskip=half}}
\makeatother
\usepackage{longtable,booktabs,array}
\usepackage{calc} % for calculating minipage widths
% Correct order of tables after \paragraph or \subparagraph
\usepackage{etoolbox}
\makeatletter
\patchcmd\longtable{\par}{\if@noskipsec\mbox{}\fi\par}{}{}
\makeatother
% Allow footnotes in longtable head/foot
\IfFileExists{footnotehyper.sty}{\usepackage{footnotehyper}}{\usepackage{footnote}}
\makesavenoteenv{longtable}
\usepackage{graphicx}
\makeatletter
\newsavebox\pandoc@box
\newcommand*\pandocbounded[1]{% scales image to fit in text height/width
  \sbox\pandoc@box{#1}%
  \Gscale@div\@tempa{\textheight}{\dimexpr\ht\pandoc@box+\dp\pandoc@box\relax}%
  \Gscale@div\@tempb{\linewidth}{\wd\pandoc@box}%
  \ifdim\@tempb\p@<\@tempa\p@\let\@tempa\@tempb\fi% select the smaller of both
  \ifdim\@tempa\p@<\p@\scalebox{\@tempa}{\usebox\pandoc@box}%
  \else\usebox{\pandoc@box}%
  \fi%
}
% Set default figure placement to htbp
\def\fps@figure{htbp}
\makeatother
\setlength{\emergencystretch}{3em} % prevent overfull lines
\providecommand{\tightlist}{%
  \setlength{\itemsep}{0pt}\setlength{\parskip}{0pt}}
\usepackage{booktabs}
\usepackage{longtable}
\usepackage{array}
\usepackage{multirow}
\usepackage{wrapfig}
\usepackage{float}
\usepackage{colortbl}
\usepackage{pdflscape}
\usepackage{tabu}
\usepackage{threeparttable}
\usepackage{threeparttablex}
\usepackage[normalem]{ulem}
\usepackage{makecell}
\usepackage{xcolor}
\usepackage{bookmark}
\IfFileExists{xurl.sty}{\usepackage{xurl}}{} % add URL line breaks if available
\urlstyle{same}
\hypersetup{
  hidelinks,
  pdfcreator={LaTeX via pandoc}}

\author{}
\date{\vspace{-2.5em}}

\begin{document}

\begin{longtable}[]{@{}l@{}}
\toprule\noalign{}
\endhead
\bottomrule\noalign{}
\endlastfoot
tle: ``Technical Report'' \\
thor: ``Emily Wu, Starr Mark, Sophie Kemprecos, Hannah Vollen'' \\
te: ``2026-04-15'' \\
tput: pdf\_document \\
\end{longtable}

\section{Background and Resarch
Questions}\label{background-and-resarch-questions}

Disciplinary infractions during incarceration reflect both individual
inmate wellbeing and broader institutional safety. This report analyzes
prisoner outcomes during their prison sentences and examines how a range
of factors influence the infractions during incarceration. This analysis
is designed to give prison administrators and policymakers a clearer
picture of where targeted investment in mental health services could
have the greatest impact. Using a current metric of wellness through
violations has been shown to be predictive of how prisoner behavior may
evolve after their sentence is over. Rule violations during
incarceration offer insight into inmate wellbeing and institutional
stability that earlier outcome studies have largely overlooked.

This analysis is intended to help identify areas where inmates are at
their most vulnerable so that limited prison resources can be used in
the most efficient ways. By using rule violations, we hope to push for
earlier intervention in mental health treatment and identifying at-risk
individuals. This is a shift from other analyses which often look at
post-sentencing outcomes to determine wellness.

Understanding who is most at risk of violations and the role mental
health plays in that risk is the missing piece that this project
addresses. There is discussion surrounding mental health in the general
population that rarely considers the prison population. The rates of
mental health issues are much higher for incarcerated people and those
with any kind of a prison record. While treatment opportunities do exist
for prisoners, the offerings typically do not have a large variety or a
high level of quality. In this report we assess and communicate findings
about treatment at a broad level, as there was no way to identify what
type of treatment a prisoner had received. The result is a broad base
look at the pervasiveness of mental health issues and how treatment is
interacting with other factors to change rule violation outcomes.

In the first phase of our analysis, we sought to identify which
variables best distinguish high and low predicted infraction risk,
investigating potentially influential demographic, criminal justice, and
socioeconomic characteristics (gender, years of education, offense type,
number of children, etc.), and specifically whether inmates received
mental health treatment (MHT) post-admission.

In the second phase of our analysis, we sought to reveal which
demographic and clinical factors drive higher infraction risk and
identify potential gaps in prison mental health services that could lead
to policy change(s).

Statistically, the research questions are important because they test
which variables are most useful for infraction risk prediction or
explanation in the presence of several predictors at a time. They also
help to distinguish prediction from explanation, which is important
because a variable can statistically be associated with risk without
being a direct cause.

Substantively, the research questions are important because infractions
impact prison safety, management, expense, and inmate outcomes.
Investigating demographic, criminal justice, and mental health treatment
factors (among others) can identify which groups are at higher risk and
whether prison mental health services are reaching those who need them
most.

\section{Data Description}\label{data-description}

We used data from the \emph{Survey of Prison Inmates (SPI), United
States, 2016 (ICPSR 37692) Data Set for Public-Use}. Data was collected
through in-person interviews with prisoners incarcerated in state
correctional facilities, as well as from inmate records, intake forms,
surveys, questionnaires, screenings, etc.

The Principal Investigator of this study was the Bureau of Justice
Statistics Conducted by the Bureau of Justice Statistics, this was
conducted such that policymakers could address shortcomings of
correctional institutions. Sample selection was done carefully, with a
two stage probability sample, where about 20,000 inmates were chosen to
represent about 1.5 Million inmates within the prison system. Data was
collected using Computer-Assisted Personal Interviewing to assist with
human and data entry error with a relatively high response rate (70\%).

Questions regarding personal feelings on past wrongdoings may lead to
response biases as inmates may feel they want to lie to ``look good.''
Inmates were assured that answers were completely anonymous and that
their identity was protected, to reduce bias.

Our model is composed of variables related to mental health, treatment,
alcohol use, and a host of controls that are unrelated to treatment. The
variable that outlines the number of diagnoses an inmate has is highly
relevant to our analysis, as it serves as a proxy for severity of
someone's mental condition. People who are diagnosed with more mental
health issues are expected to struggle more in and outside of prison. On
top of that, the mental health treatment status of an inmate is
important to this work as it is one of the focal parameters. It
indicates if a person has received mental health treatment during their
incarceration period or not. Additionally we are able to observe if the
inmate has had an alcohol abuse problem or if they have ever
participated in an alcohol treatment program.

In terms of the variables outside of treatment and mental health,
education is one of the most important variables in our analysis. In the
exploratory data analysis of the SPI, 78 percent of prisoners had a high
school diploma or less. In the second stage model, we find education to
be one of the strongest predictors of rule violations. Additionally, we
find that the number of times a person is arrested is highly
significant. This is related to the idea of repeated violations of legal
rules that likely transfers over to prison-time violations.

Many of the variables were suppressed due to privacy reasons and
therefore had to be removed because they were unusable. Additionally,
all categorical variables had extra descriptions, i.e.~``(1) 1 = Yes''
as an observation, and therefore must be flattened. Furthermore, many
observations were missing for various reasons with codes to specify why,
but for our purposes, we flattened them all to ``0'' or NA as the reason
for missing was not relevant to this project.

\subsection{Exploratory Data Analysis}\label{exploratory-data-analysis}

We conducted exploratory data analysis to further our understanding of
the variables in our data set. Using data visualizations of the
distribution of variables and correlation matrices allowed us to see the
range of outcomes for prison inmates across variables like education,
mental health treatment, and gender. Identifying if there are
differences in sub groups was informative to seeing the variation across
risk groups of inmates for prison violations.

\pandocbounded{\includegraphics[keepaspectratio]{Technical_files/figure-latex/unnamed-chunk-2-1.pdf}}

The above figure is a boxplot of female versus male inmates from the
sample. In exploration of what kind of variation there might be between
inmate subgroups, we assessed gender as the literature had often noted
strong differences in inmate behavior based on gender. The female prison
population has a much more concentrated number of violations around zero
in comparison to the male population. The male inmates share a similar
median around one rule violation, but have a much larger spread of data
compared to the female inmates. This indicated that gender would be an
important demographic variable to include in our model.

\pandocbounded{\includegraphics[keepaspectratio]{Technical_files/figure-latex/unnamed-chunk-3-1.pdf}}

The data includes variables that detail if a prisoner has ever received
a diagnosis of bipolar disorder, depression, schizophrenia, anxiety,
PTSD, and a host of other mental health issues. We created a variable
that summarizes the number of diagnoses based on the existing diagnoses
an inmate has. The above boxplot shows a moderate infraction count for
those with two or fewer diagnoses, but a noticeable jump in spread
occurs at the third diagnoses. There is a noticeable increase in the
count of infractions after the third diagnosis, a trend that continues
up until the sixth mental health diagnoses.

\pandocbounded{\includegraphics[keepaspectratio]{Technical_files/figure-latex/unnamed-chunk-4-1.pdf}}

The offense type also had an interesting distribution. Intuition
suggested that the most violent offenses would also have the highest
shares of rule violaters compared to non-violators. This holds true, but
three other offense types have similarly high shares of rule violators:
drug, property, and public offenses. This identifies a meaningful
difference in the distribution of violators motivating why offense type
is included in our later regressions.

\section{Methods}\label{methods}

Our analysis uses a two-stage modeling approach to answer our research
questions. In the first stage, we implement a logistic regression model
to classify inmates into ``high'' and ``low'' predicted risk groups
based on whether they are at risk of committing rule violations. This
model uses demographic, criminal justice, and mental health-related
variables to estimate the probability that an individual falls into the
at-risk category. By transforming the outcome into a binary indicator,
this stage allows us to identify which characteristics are most
predictive of elevated violation risk and to create a clear distinction
between higher and lower risk individuals.

In the second stage of our analysis, we used a negative binomial
regression model to the subset of inmates classified as high predicted
risk. In this stage, the outcome variable is the number of rule
violations, which is a count variable with a wide range of possible
values. The negative binomial model allows us to estimate how individual
characteristics are associated with variation in the number of
infractions within this higher-risk group. This two-stage approach
allows us to first identify who is at risk and then more precisely model
the intensity of rule violations among those individuals.

The negative binomial regression is the best fit for our second stage
because our outcome variable is a count of rule violations that shows
clear overdispersion. When we initially used a Poisson model, we found
that the model had a much higher AIC and did not appropriately account
for the variance in the data, which was much larger than the mean.
Because the Poisson model assumes that the mean and variance are equal,
this led to a poor fit for our data. The negative binomial model relaxes
this assumption and allows the variance to exceed the mean, making it
more appropriate for modeling the distribution of infractions in our
sample.

Inspection of the upper tail of the distribution revealed several
extreme values (e.g., 50, 54, 60, 63, 70, 78, 90, 100, 120, 175, 250)
that were substantially larger than the rest of the data. These values
were sparsely distributed and exhibited large jumps, suggesting they may
reflect institutional recording practices or aggregation of infractions
rather than comparable counts of individual violations. To reduce the
influence of these extreme observations on model estimation, values
above 50 infractions were excluded from the analysis.

Additionally, the structure of our data includes a high number of
observations at or near zero, along with a smaller number of individuals
with a large number of violations. The negative binomial model is better
suited to handle this type of dispersion and skewness, allowing for more
accurate and stable coefficient estimates. This is particularly
important given the wide range of possible outcomes for the number of
infractions.

We also considered alternative modeling approaches. A Poisson regression
was initially explored, but ultimately rejected due to its poor model
fit and inability to handle overdispersion. We also considered using a
single-stage model rather than a two-stage approach, but this would have
forced us to model both the probability of being at risk and the number
of infractions simultaneously, which could reduce interpretability and
model accuracy given the zero-inflated nature of the data. By separating
the process into two stages, we are able to more clearly distinguish
between the factors that predict risk and the factors that influence the
number of violations among high-risk individuals.

The first stage of the analysis using a logistic regression. The
logistic regression relies on the assumption of there being a linear
relationship between the outcome variable, if a prisoner is at risk, and
the independent variables. These independent variables are our
regressors of age, gender, race, military status, offense type, mental
health treatment, and diagnoses. In order to use the logistics
regression to evaluate the relationship, the sample size must be large.
In our sample we begin with over 20,000 inmates. The regression uses
15,000 inmate records for data which are complete across the variables
used which is a very large model sampling size.

In a logistic regression the outcome variable has to be binary, with a 0
or 1 outcome indicator. To achieve this we created the ``at risk''
variable. This variable captures if prisoners are expected to commit
rule violations, sorted into a yes or no grouping associated with a 1
and 0. We also do not believe that the observations are correlated at
all. Each inmate should be an independent, randomly sampled unit from
across the state prison population. There are no federal prisoners
represented in our data, but this sample of people should be
representative of the larger prison population and not lead to biased
results from a dependence in observation responses. To avoid
multicollinearity in the logistic model, we reviewed exploratory data
analysis to evaluate if any variables had a strong relationship with
each other. We found that bipolar disorder and a depression diagnosis
were highly correlated with each other. The depression variable was
noticeably correlated with the other diagnoses, this may be due to a
high presence of depression found across those with other mental health
issues.

One assumption the negative binomial has is independence of the
observations. This is mainly relevant to those working in longitudinal
panel data sets. This is not a concern for our model as we have snapshot
data of prison inmates limited to one year compared to a multi-year
tracking panel. Negative binomial models perform best in large sample
sizes. Small sample sizes with highly dispersed results will lead to
inaccurate coefficient estimates. Given our data set has more than
10,000 observations, we feel the assumption of large sample sizes is
met.

The negative binomial regression also can not be used if there is a high
presence of multicollinearity. We performed exploratory data analysis to
look at inter-variable correlation rates. To ensure that no two
variables were accounting for each other we checked a correlation matrix
and heatmap of the explanatory variables and did not include those with
correlations above \textbar0.7\textbar. This was done by evaluating the
different explanatory plots from the logistic regression and then
checked the VIF levels to ensure they were below 10.

\section{Results}\label{results}

Our Final First Stage Model is as follows:

\(logit(P(At Risk = 1)) = \beta_0+\beta_1Years of Education+\beta_2ADD/ADHD+\beta_3Gender:Male+\beta_4Offense:Other+\beta_5Offense:Property+\beta_6Offense:Public Order+\beta_7Offense:Unknown+\beta_8Offense:Violent+\beta_9Children+\beta_{10}Times Arrested+\beta_{11}Mental Health Treatment+\beta_{12}Alcoholism\)

The first stage of analysis employed logistic regression, a binomial
family logit link model. The reported coefficients, \(\beta\), represent
the change in log-odds of being classified as ``at risk'' for every
one-unit increase in the predictor variable. These estimates are modeled
according to the logit link function:

\(logit(p) = \ln(\frac{p}{1-p}) = \beta_0+\beta_1x_1...\beta_kx_k\)

where \(p\) is the probability of risk classification

After conducting a Likelihood-Ratio Test on our model, we find that our
final model has a deviance score (null vs residual) of 4685 on 12
degrees of freedom and a p-value less than 0.001. This model was fit on
\(n=15,450\) observations with 1708 dropped/deleted observations. Over
90\% are due to alcoholism.

This indicates our full model significantly outperformed the null model
and the combination of clinical and demographic factors provides
substantial predictive power for inmate infraction risk. Overall, our
model had a 73\% classification accuracy (classification of inmates who
had infractions or not) and with an AUC of 0.797, it demonstrates a good
amount of discriminative ability, correctly distinguishing between at
risk and non-at risk individuals about 80\% of the time when presented
with a random pair of one positive and one negative case.

\begin{ThreePartTable}
\begin{TableNotes}
\item \textit{Note: } 
\item Overall Accuracy: 73.3\% | Sensitivity: 75.3\% | Specificity: 71.6\% | P-Value: < .001 | No Information Rate (NIR): 53.9\%
\end{TableNotes}
\begin{longtable}[t]{lcc}
\caption{\label{tab:unnamed-chunk-5}Stage 1 Confusion Matrix}\\
\toprule
Classification & Observed: No Infraction (0) & Observed: At Risk (1)\\
\midrule
Predicted: Low Risk (0) & 5,361 & 2,365\\
Predicted: High Risk (1) & 1,755 & 5,969\\
\bottomrule
\insertTableNotes
\end{longtable}
\end{ThreePartTable}

\pandocbounded{\includegraphics[keepaspectratio]{Technical_files/figure-latex/unnamed-chunk-5-1.pdf}}

The model diagnostics state that the most significant predictors are
ADD/ADHD, Male Gender, Times Arrested, and Mental Health Treatment. Out
of these, the predictors with the largest effect size are ADD/ADHD with
1.33, Male Gender with -0.938, and Mental Health Treatment with 2.201.

\begin{longtable}[]{@{}
  >{\raggedright\arraybackslash}p{(\linewidth - 8\tabcolsep) * \real{0.3380}}
  >{\raggedleft\arraybackslash}p{(\linewidth - 8\tabcolsep) * \real{0.1549}}
  >{\raggedleft\arraybackslash}p{(\linewidth - 8\tabcolsep) * \real{0.1408}}
  >{\raggedleft\arraybackslash}p{(\linewidth - 8\tabcolsep) * \real{0.1408}}
  >{\raggedright\arraybackslash}p{(\linewidth - 8\tabcolsep) * \real{0.2254}}@{}}
\caption{Significant Predictors}\tabularnewline
\toprule\noalign{}
\begin{minipage}[b]{\linewidth}\raggedright
Variable
\end{minipage} & \begin{minipage}[b]{\linewidth}\raggedleft
Estimate
\end{minipage} & \begin{minipage}[b]{\linewidth}\raggedleft
OR
\end{minipage} & \begin{minipage}[b]{\linewidth}\raggedleft
Z\_Value
\end{minipage} & \begin{minipage}[b]{\linewidth}\raggedright
95\% Confint IRR
\end{minipage} \\
\midrule\noalign{}
\endfirsthead
\toprule\noalign{}
\begin{minipage}[b]{\linewidth}\raggedright
Variable
\end{minipage} & \begin{minipage}[b]{\linewidth}\raggedleft
Estimate
\end{minipage} & \begin{minipage}[b]{\linewidth}\raggedleft
OR
\end{minipage} & \begin{minipage}[b]{\linewidth}\raggedleft
Z\_Value
\end{minipage} & \begin{minipage}[b]{\linewidth}\raggedright
95\% Confint IRR
\end{minipage} \\
\midrule\noalign{}
\endhead
\bottomrule\noalign{}
\endlastfoot
ADD/ADHD & 1.3389959 & 3.8152109 & 28.22076 & (3.48 - 4.19) \\
Male Gender & -0.9386538 & 0.3911541 & -19.65422 & (0.36 - 0.43) \\
Mental Health Treatment & 2.2016663 & 9.0400643 & 43.74620 & (8.2 -
9.98) \\
\end{longtable}

\newpage

The final Second Stage Model is as follows:

\(ln(Infractions) = \beta_0+\beta_1Race:Black+\beta_2Race:Hispanic+\beta_3Race:Multiracial+\beta_4Race:Other+\beta_5Race:White+\beta_6NumberofDiagnoses+\beta_7ArmedForces+\beta_8TimesArrested+\beta_9YearsofEducation+\beta_{10}ADD/ADHD+\beta_{11}Homeless+\beta_{12}Children+\beta_{13}MentalHealthTreatment+\beta_{14}Alcoholism+\beta_{15}AlcoholTreatment+\beta_{16}Gender:Male+\beta_{17}+ \beta_{18}Offense:Other+\beta_{19}Offense:Property+\beta_{20}Offense:PublicOrder+\beta_{21}Offense:Unknown+\beta_{22}Offense:Violent+\beta_{23}NumberofDiagnoses*MentalHealthTreatment+\beta_{24}Alcoholism*Alcohol Treatment\)

The second stage utilized negative binomial regression with a log link
to analyze the infraction frequency within the high-risk subgroup. The
coefficients are reported on the log-count scale. Due to the model using
the natural log link, the relationship is such:

\(\ln(\mu) =\beta_0+\beta_1x_1...\beta_kx_k\)

Comparing the 2 models using their goodness-of-fit statistics, it is
seen that a negative binomial model is more favorable with a lower AIC
and Overdispersion statistic \((\phi)\). The negative binomial model was
fitted on the subpopulation of \(n = 7,619\) of those rated as ``High
Predicted Risk.'' This was determined by our previous logistic model of
those who had a higher predicted probability of having an infraction
than the median value of the data set, which was 0.497. 105 observations
were not counted, over 90\% were deleted due to Race being unavailable.

\begin{longtable}[]{@{}lrrrr@{}}
\caption{Model Comparison}\tabularnewline
\toprule\noalign{}
Model & AIC & Deviance & df & Overdispersion \\
\midrule\noalign{}
\endfirsthead
\toprule\noalign{}
Model & AIC & Deviance & df & Overdispersion \\
\midrule\noalign{}
\endhead
\bottomrule\noalign{}
\endlastfoot
Poisson & 35484.49 & 2373.456 & 23 & 8.3143945 \\
Negative Binomial & 19761.59 & 385.641 & 23 & 0.2250685 \\
\end{longtable}

In our model, the most significant predictors are ADD/ADHD, Times
Arrested, Children, Mental Health Treatment, Alcoholism, and an Unknown
Offense Type. The variables with the largest effect size were ADD/ADHD
with a 0.605 Coefficient and Unknown Offense type with -0.513
coefficient.

\begin{longtable}[]{@{}
  >{\raggedright\arraybackslash}p{(\linewidth - 8\tabcolsep) * \real{0.3478}}
  >{\raggedleft\arraybackslash}p{(\linewidth - 8\tabcolsep) * \real{0.1594}}
  >{\raggedleft\arraybackslash}p{(\linewidth - 8\tabcolsep) * \real{0.1449}}
  >{\raggedright\arraybackslash}p{(\linewidth - 8\tabcolsep) * \real{0.1159}}
  >{\raggedright\arraybackslash}p{(\linewidth - 8\tabcolsep) * \real{0.2319}}@{}}
\caption{Model Statistics}\tabularnewline
\toprule\noalign{}
\begin{minipage}[b]{\linewidth}\raggedright
Variable
\end{minipage} & \begin{minipage}[b]{\linewidth}\raggedleft
Estimate
\end{minipage} & \begin{minipage}[b]{\linewidth}\raggedleft
Std.Error
\end{minipage} & \begin{minipage}[b]{\linewidth}\raggedright
P.Value
\end{minipage} & \begin{minipage}[b]{\linewidth}\raggedright
95\% Confint IRR
\end{minipage} \\
\midrule\noalign{}
\endfirsthead
\toprule\noalign{}
\begin{minipage}[b]{\linewidth}\raggedright
Variable
\end{minipage} & \begin{minipage}[b]{\linewidth}\raggedleft
Estimate
\end{minipage} & \begin{minipage}[b]{\linewidth}\raggedleft
Std.Error
\end{minipage} & \begin{minipage}[b]{\linewidth}\raggedright
P.Value
\end{minipage} & \begin{minipage}[b]{\linewidth}\raggedright
95\% Confint IRR
\end{minipage} \\
\midrule\noalign{}
\endhead
\bottomrule\noalign{}
\endlastfoot
ADD/ADHD & 0.6047363 & 0.0637578 & \textless0.0001 & (1.616 - 2.074) \\
Times Arrested & 0.0079125 & 0.0011492 & \textless0.0001 & (1.006 -
1.01) \\
Children & -0.3953537 & 0.0588810 & \textless0.0001 & (0.6 - 0.756) \\
Mental Health Treatment & 0.3248295 & 0.0936386 & 0.00052 & (1.152 -
1.663) \\
Alcoholism & 0.2941803 & 0.0815311 & 0.00031 & (1.144 - 1.575) \\
Offense Type: Unknown & -0.5126795 & 0.1334874 & 0.00012 & (0.461 -
0.778) \\
\end{longtable}

\pandocbounded{\includegraphics[keepaspectratio]{Technical_files/figure-latex/unnamed-chunk-9-1.pdf}}

Diagnostic plots suggest that the negative binomial model reproduces the
broad distributional form of the outcome, but does not fully capture the
allocation of probability among the smallest counts. The largest
prediction errors occur at 0, 1, and 2 infractions, while fit improves
for higher counts. This indicates that the model is adequate for the
general shape of the count distribution but less accurate in modeling
the concentration of low-count outcomes.

\pandocbounded{\includegraphics[keepaspectratio]{Technical_files/figure-latex/unnamed-chunk-10-1.pdf}}

\section{Interpretation}\label{interpretation}

Our Stage One results indicate that high‑risk classification in this
sample is driven primarily by clinical factors and criminal justice
history rather than socioeconomic variables. An ADD/ADHD diagnosis and
receipt of post‑admission mental health treatment are the strongest
predictors of being labeled high risk. The coefficient for ADD/ADHD
(\(\approx\) 1.34) corresponds to an odds ratio of about 3.8, meaning
individuals with ADD/ADHD have nearly four times the odds of being
classified as high risk compared to those without this diagnosis,
holding other covariates constant. The coefficient for mental health
treatment (\(\approx\) 2.02; odds ratio \(\approx\) 7.5) indicates that
inmates who receive mental health treatment post‑admission have more
than seven times the odds of being classified as high risk relative to
those who do not receive treatment. Criminal justice variables also play
an important role: each additional prior arrest is associated with a
modest but statistically significant increase in the odds of high‑risk
classification (about 1--2\% per arrest, cumulating over many arrests),
and each additional child is associated with roughly a 9\% increase in
the odds of high‑risk status. Property and public‑order offenses exhibit
higher odds of high‑risk classification than the reference offense type,
while years of education are not significantly related to risk status.
Notably, the coefficient for male gender is negative and large in
magnitude (odds ratio \(\approx\) 0.39), implying that, conditional on
diagnoses, treatment, offense type, and criminal history, males in our
sample are less likely than the reference gender category to be
classified as high risk. In Stage 2, we restricted the sample to inmates
classified as ``High Predicted Risk'' and modeled the count of
institutional infractions using a negative binomial regression. Within
this high‑risk subgroup, both structural disadvantage and clinical
burden are central determinants of infraction counts. Each additional
prior arrest and each additional mental health diagnosis are associated
with approximately a 7\% increase in the expected number of infractions
(incidence‑rate ratios \(\approx\) 1.07), suggesting that deeper
criminal justice involvement and greater psychiatric complexity are
linked to more institutional misconduct. In contrast, years of education
are protective once we condition on being high risk: each additional
year of schooling is associated with about a 7\% reduction in expected
infractions (IRR \(\approx\) 0.93), implying that more educated
high‑risk inmates accumulate fewer infractions. Alcoholism is associated
with substantially higher infraction counts (IRR \(\approx\) 1.34),
whereas ADD/ADHD is no longer a significant predictor once we restrict
to the high‑risk group. Demographically, male inmates have higher
infraction counts than the reference gender group (IRR \(\approx\) 1.5),
and Black inmates exhibit markedly higher infraction counts than the
reference race category (IRR on the order of 2.8), indicating large
racial disparities in disciplinary outcomes among high‑risk inmates even
after adjustment for clinical and criminal‑history covariates. Offense
type is less influential in Stage 2; most categories are not
significantly associated with infraction counts, with the exception of
an ``unknown'' offense type that is associated with fewer infractions.

Treatment variables display different patterns. Mental health treatment
remains positively associated with infraction counts among high‑risk
inmates. The coefficient of roughly 0.31 (IRR \(\approx\) 1.37)
indicates that treated high‑risk inmates have about 37\% more
infractions than untreated high‑risk inmates, conditional on diagnoses
and other controls. An interaction between the number of diagnoses and
mental health treatment is small and non-significant, suggesting that
mental health treatment does not materially alter the relationship
between diagnostic burden and infractions in our data. By contrast,
alcohol/drug treatment shows a clinically meaningful interaction with
alcoholism. While the main effect of alcohol/drug treatment is
negligible overall, the Alcoholism × Alcohol/Drug Treatment interaction
term is negative and statistically significant (IRR \(\approx\) 0.76),
indicating that among inmates with alcoholism, those who receive
alcohol/drug treatment have approximately 24\% fewer expected
infractions than alcoholic inmates who do not receive such treatment.

Our first research question asked which variables best distinguish
inmates with high versus low predicted infraction risk, with particular
interest in demographic, criminal justice, and socioeconomic
characteristics, and in whether inmates received mental health treatment
post‑admission. The Stage 1 logistic regression results provide a clear
answer. High‑risk classification is primarily shaped by clinical and
justice‑system variables---ADD/ADHD, alcoholism, receipt of mental
health treatment, number of prior arrests, and offense type---rather
than by socioeconomic status as proxied by years of education. Mental
health treatment is strongly and positively associated with high‑risk
status, implying that mental health services are being targeted toward
inmates whom the system already identifies as behaviorally or clinically
high risk. Gender and family structure also contribute: males are less
likely to be flagged as high risk once other factors are controlled, and
inmates with more children are more likely to be classified as high
risk. These findings address the first stage of our questions by
demonstrating which factors the model uses to separate high‑ and
low‑risk inmates and how mental health treatment is distributed across
those groups. Our second research question focused on the determinants
of infraction severity within the high‑risk group and the extent to
which patterns in mental health and substance‑use treatment reveal
potential gaps in services. The Stage 2 negative binomial model speaks
directly to this. Within the high‑risk subgroup, higher infraction
counts are associated with more prior arrests, more mental health
diagnoses, alcoholism, lower levels of education, male gender, and Black
race. Thus, our results identify a set of demographic and clinical
factors---especially cumulative diagnostic burden and substance‑use
disorders---that drive higher infraction risk among already high‑risk
inmates. The treatment variables further inform our understanding of
potential gaps. Mental health treatment is highly concentrated among
inmates with greater clinical and behavioral severity and is positively
associated with infractions, which is consistent with severity‑driven
selection into services rather than clear behavioral improvement. In
contrast, the interaction between alcoholism and alcohol/drug treatment
suggests a potentially beneficial effect of substance‑use services among
those with alcoholism, as treated alcoholic inmates have fewer
infractions than untreated alcoholic inmates. These patterns address our
second question by highlighting both who is most at risk within the
high‑risk population and where current mental health and substance‑use
services may or may not be reducing that risk.

Based on these models, we conclude that high‑risk classification in this
facility is closely tied to specific clinical conditions and
criminal‑history indicators. Inmates with ADD/ADHD, alcoholism, more
arrests, more children, and certain offense types are substantially more
likely to be flagged as high risk, whereas years of education do not
play a meaningful role at the classification stage. Among those already
classified as high risk, the intensity of institutional misconduct is
shaped by cumulative mental health burden, substance‑use status,
criminal justice history, and structural disadvantage. In particular,
additional diagnoses and prior arrests are associated with more
infractions, alcoholism is a strong risk factor, and more years of
education are associated with fewer infractions. The higher infraction
rates observed among male and Black high‑risk inmates point to important
gender and racial disparities in disciplinary outcomes that warrant
further investigation.

With respect to services, our findings suggest that mental health
treatment is being targeted toward the highest‑risk inmates but, in
these observational models, is not associated with lower infraction
counts; the positive association likely reflects confounding by severity
rather than a harmful effect of treatment. By contrast, alcohol/drug
treatment appears to be behaviorally beneficial within the subgroup of
inmates with alcoholism, as indicated by the significant interaction
term. Taken together, these results support policy recommendations to
intensify screening and treatment for alcohol use disorders among
high‑risk inmates, to expand or strengthen education‑related programming
for this population, to re‑examine the design and delivery of mental
health services with respect to behavioral outcomes, and to scrutinize
racial and gender disparities in disciplinary practices and infraction
recording.

\section{Limitations}\label{limitations}

Treatment participation is not randomly assigned, and individuals
receiving treatment may differ systematically from those who do not.
This introduces potential selection bias that may confound estimated
treatment effects. As a result, estimated treatment effects should be
interpreted as associative rather than causal, and may reflect
underlying differences in inmate characteristics rather than the effect
of treatment itself.

Additionally, several variables, including mental health diagnoses and
substance use, are self-reported and may be subject to measurement error
or misclassification. This is best shown in the ``Alcoholism'' Variable.
A substantial portion of missing data in the SPI is due to this and
therefore affected the first-stage logistic model. If missingness in
this variable is systematically related to unobserved characteristics,
then the resulting sample may not be representative of the full inmate
population. Many of the drug variables were missing and ultimately, had
to be taken out of the model due to an excessive amount of observations
being unusable. This introduces the potential for selection bias and may
lead to biased coefficient estimates, particularly for variables
correlated with substance use and disciplinary outcomes.

The second stage model entirely relies on the logistic model to be
strong. The initial logistic regression must be a resilient and robust
predictor of if a prisoner is in the at risk for high prison violations
or not. If the regression turns out to be a weak predictor, then any
results from the second stage will be highly criticized for their
dependence on a weak dependent variable. The validity of the second
stage depends on the predictive performance of the first-stage model. If
the classification model performs poorly, misclassification of
individuals into the high-risk group may introduce bias into the
second-stage estimates.

An important limitation of trying to predict rule violations is that we
do not know exactly when the violations occurred. The relevancy of
timing is important to highlight, but the data did not allow for highly
granular details on when and what kind of rule violations were
committed. There may have been important information on the severity of
the infractions that we are missing in our model or outcomes. Given the
limits of our data themself we felt that the reduced form model will
still yield general trends that are important for prison populations but
won't identify what kind of rule violations are occurring. Additionally
there is data on location and prison identification that was not
available in the publicly accessible data. If we had access to location
then we could control for state differences in prison systems and try to
limit the scope of the population to a more defined population. This
leaves us with very high-level takeaways for a general population of
state prison inmates that informs policy makers and stakeholders about
general trends rather than specific insights, limiting causal
interpretation. These suggest that the results should be interpreted as
descriptive and predictive rather than causal.

For the second stage, we are using a negative binomial regression. The
negative binomial regression models the expected count of infractions
while accounting for overdispersion in the data. With rule violations as
the outcome variable, we see potential outcomes that range from one to
100. In order to model such a large dispersion of outcomes the negative
binomial regression has weaknesses that are key to highlighting in our
study.

If there are more zero values than the model expects, the results for
beta coefficients will be biased. Given the concentration of the outcome
variable's response around zero there may be bias in our estimates. The
high concentration of zero outcomes may indicate potential
zero-inflation, which is not explicitly modeled. The two-stage framework
partially addresses this by conditioning on high-risk individuals, but
it does not fully resolve distributional misspecification. Another
limitation of the model is that it is influenced by outliers and
skewness. To mitigate the influence of extreme values, observations
exceeding 50 infractions were excluded. While these values appeared to
be qualitatively different from the majority of the data, their removal
may exclude valid high-frequency cases and limit the model's ability to
capture behavior among the most extreme offenders.

\section{Conclusion}\label{conclusion}

Taken all together, the 2 stage analysis shows that the prison
disciplinary risk is shaped primarily by clinical burden and criminal
justice history, and not simply by broad socioeconomic background. In
the first stage, inmates were most likely to be classified as high-risk
when they had markers of greater psychiatric or behavioral severity,
such as ADD/ADHD, alcoholism, prior arrests, and post-admission mental
health treatment. This suggests that the classification process is
strongly tied to observable mental health needs and prior system
involvement.

In the second stage, once attention is restricted to inmates already
identified as high risk, the intensity of misconduct is driven by a
somewhat different set of factors: more prior arrests, greater
cumulative mental health disorder diagnoses, alcoholism, lower levels of
education, and mental health treatment are all associated with higher
expected infraction counts, indicating that both clinical complexity and
structural disadvantage continue to matter within the highest risk
subgroup. Interestingly, the effect of male gender becomes not
significant when in the presence of other variables in the second stage
model. Although males make up the majority of the overall sample
(approximately 75\%), they represent a smaller share of the high-risk
group (about 60\%). Conversely, females comprise a larger proportion of
the high-risk group (approximately 40\%) compared to their
representation in the full sample (about 25\%), indicating that females
are disproportionately represented among individuals classified as high
risk. Mental health treatment remains positively associated with
disciplinary outcomes, but in the context of these observational data
this pattern is most consistent with severity-based selection into
treatment rather than evidence that treatment worsens behavior. By
contrast, alcohol/drug treatment appears to be more promising for the
subgroup with alcoholism, as the negative and significant interaction
suggests that inmates with alcoholism who receive substance-use
treatment accumulate fewer infractions than comparable inmates who do
not receive such services.

To reiterate, treatment effects cannot be understood apart from inmate
severity. Mental health services appear to be concentrated among those
with greatest behavioral and psychiatric needs, while substance-use
treatment shows clearer evidence of association with reduced misconduct
in the population for whom it is most directly targeted. The disparate
behavioral outcomes between mental health treatment and substance use
programming may be attributed to differences in sociocultural stigma and
voluntary engagement. While mental health services in correctional
facilities are often crisis-oriented and carry a clinical `labeling'
burden, alcohol and drug treatments frequently utilize peer-support and
community-based models. This may facilitate higher-quality engagement
and a greater internal motivation for behavioral change, potentially
explaining the significant interaction effect observed among inmates
with alcoholism. Without the availability of temporal data, causal
inference cannot be employed in our study. More broadly, our findings
suggest that the inmates most vulnerable to repeated disciplinary
problems are those carrying multiple forms of disadvantage at once. The
positive correlation between mental health treatment and infractions
likely reflects severity-driven selection. Prior justice-system contact,
psychiatric burden, substance-use disorder, and, within the high-risk
group, marked racial disparities in disciplinary outcomes. As a result,
the strongest policy implications of this study are to intensify early
screening for alcohol-use disorder, strengthen targeted substance-use
treatment for high-risk inmates, expand educational programming as a
potential protective factor, and more closely examine whether mental
health services are sufficiently aligned with behavioral outcomes.

Future work would benefit from data that include the timing of treatment
and infractions, institutional identifiers, and more detailed measures
of offense severity, since those additions would make it possible to
distinguish more clearly between treatment targeting, treatment
effectiveness, and differences in disciplinary practices across
facilities. The SPI data includes an extensive list of variables as the
survey was very comprehensive. During the exploratory data analysis
phase, post-incarceration socioeconomic indicators, specifically
employment status and income, were considered as potential outcome
measures for mental health treatment efficacy. However, these variables
were excluded from the final analysis as they are outside the data
collection scope of the Survey of Prison Inmates. Perhaps this could be
an angle for data collection to find relationships between efficacy
outside of prison.

\newpage

\section{Appendix A: Data Dictionary}\label{appendix-a-data-dictionary}

\begingroup\fontsize{7}{9}\selectfont

\begin{longtable}[t]{>{\raggedright\arraybackslash}p{2.5cm}>{\raggedright\arraybackslash}p{1.5cm}>{\raggedright\arraybackslash}p{1.5cm}>{\raggedright\arraybackslash}p{4cm}>{\raggedright\arraybackslash}p{3.5cm}>{\raggedright\arraybackslash}p{3cm}}
\caption{\label{tab:unnamed-chunk-11}Codebook: Variable Descriptions and Coding}\\
\toprule
Variable Name & Original Code & Type & Description & Coding & Source\\
\midrule
\endfirsthead
\caption[]{Codebook: Variable Descriptions and Coding \textit{(continued)}}\\
\toprule
Variable Name & Original Code & Type & Description & Coding & Source\\
\midrule
\endhead

\endfoot
\bottomrule
\endlastfoot
\addlinespace[0.3em]
\multicolumn{6}{l}{\textbf{Derived Variables}}\\
\hspace{1em}\cellcolor{gray!10}{at\_risk} & \cellcolor{gray!10}{Derived} & \cellcolor{gray!10}{Binary (0/1)} & \cellcolor{gray!10}{Binary indicator from median split of logistic risk probability} & \cellcolor{gray!10}{1 = Above median predicted risk probability, 0 = Below median} & \cellcolor{gray!10}{Derived from Stage 1 logistic model predictions}\\
\hspace{1em}num\_diagnoses & Derived & Numeric & Count of 6 specific mental health diagnoses & 0-6 (sum of: Bipolar, Depression, Schizophrenia, PTSD, Anxiety, and Personality Disorder) & Derived from psychiatric screening\\
\hspace{1em}\cellcolor{gray!10}{MHTx:num\_diagnoses} & \cellcolor{gray!10}{Interaction} & \cellcolor{gray!10}{Numeric} & \cellcolor{gray!10}{Mental Health Treatment x Diagnosis Count interaction} & \cellcolor{gray!10}{Product: Mental Health Treatment x num\_diagnoses} & \cellcolor{gray!10}{Model-derived interaction}\\
\hspace{1em}Alcoholism:DrugTx & Interaction & Numeric & Alcoholism x Substance Treatment interaction & Product: Alcoholism x Alcohol/Drug Treatment & Model-derived interaction\\
\addlinespace[0.3em]
\multicolumn{6}{l}{\textbf{Outcome \& Continuous Predictors}}\\
\hspace{1em}\cellcolor{gray!10}{Number Infractions} & \cellcolor{gray!10}{V1398} & \cellcolor{gray!10}{Numeric} & \cellcolor{gray!10}{Total count of disciplinary infractions during incarceration} & \cellcolor{gray!10}{Continuous (0-250+)} & \cellcolor{gray!10}{Prison disciplinary records}\\
\hspace{1em}Times Arrested & V0909 & Numeric & Lifetime number of prior arrests & Continuous (count, negatives → NA) & Criminal history records\\
\hspace{1em}\cellcolor{gray!10}{Years of Education} & \cellcolor{gray!10}{V0935} & \cellcolor{gray!10}{Numeric} & \cellcolor{gray!10}{Total years of formal education completed} & \cellcolor{gray!10}{Continuous (years, negatives → NA)} & \cellcolor{gray!10}{Self-reported education history}\\
\addlinespace[0.3em]
\multicolumn{6}{l}{\textbf{Categorical Predictors}}\\
\hspace{1em}Race & V1951-V1957 & Factor & Self-reported race/ethnicity at intake & White/Black/Hispanic/Other (Multiracial/Asian/Native Hawaiian/AI/Other collapsed) & Demographic intake form\\
\hspace{1em}\cellcolor{gray!10}{Gender} & \cellcolor{gray!10}{V1212} & \cellcolor{gray!10}{Factor} & \cellcolor{gray!10}{Self-described gender identity} & \cellcolor{gray!10}{Male/Female/Transgender/Do not identify as male, female or transgender} & \cellcolor{gray!10}{Demographic intake form}\\
\hspace{1em}offense\_type & V0062 & Factor & Primary instant offense category & Violent/Property/Drug/Public Order & Presentence investigation report\\
\hspace{1em}\cellcolor{gray!10}{Marital Status} & \cellcolor{gray!10}{V0022} & \cellcolor{gray!10}{Factor} & \cellcolor{gray!10}{Marital status at time of offense} & \cellcolor{gray!10}{0-5 (coded levels: Single/Married/Divorced/etc.)} & \cellcolor{gray!10}{Family status questionnaire}\\
\addlinespace[0.3em]
\multicolumn{6}{l}{\textbf{Binary Predictors}}\\
\hspace{1em}ADD/ADHD & V0942 & Binary (0/1) & Lifetime diagnosis of ADD/ADHD & 0 = No/Never, 1 = Yes (2 → 0 recoded) & Psychiatric screening questionnaire\\
\hspace{1em}\cellcolor{gray!10}{Bipolar} & \cellcolor{gray!10}{V1185} & \cellcolor{gray!10}{Binary (0/1)} & \cellcolor{gray!10}{Lifetime bipolar disorder diagnosis} & \cellcolor{gray!10}{0 = No, 1 = Yes (2 → 0 recoded)} & \cellcolor{gray!10}{Psychiatric screening}\\
\hspace{1em}Depression & V1186 & Binary (0/1) & Lifetime major depression diagnosis & 0 = No, 1 = Yes (2 → 0 recoded) & Psychiatric screening\\
\hspace{1em}\cellcolor{gray!10}{Schizophrenia} & \cellcolor{gray!10}{V1187} & \cellcolor{gray!10}{Binary (0/1)} & \cellcolor{gray!10}{Lifetime schizophrenia diagnosis} & \cellcolor{gray!10}{0 = No, 1 = Yes (2 → 0 recoded)} & \cellcolor{gray!10}{Psychiatric screening}\\
\hspace{1em}PTSD & V1188 & Binary (0/1) & Lifetime PTSD diagnosis & 0 = No, 1 = Yes (2 → 0 recoded) & Psychiatric screening\\
\hspace{1em}\cellcolor{gray!10}{Anxiety} & \cellcolor{gray!10}{V1189} & \cellcolor{gray!10}{Binary (0/1)} & \cellcolor{gray!10}{Lifetime anxiety disorder diagnosis} & \cellcolor{gray!10}{0 = No, 1 = Yes (2 → 0 recoded)} & \cellcolor{gray!10}{Psychiatric screening}\\
\hspace{1em}Personality Disorder & V1190 & Binary (0/1) & Lifetime personality disorder diagnosis & 0 = No, 1 = Yes (2 → 0 recoded) & Psychiatric screening\\
\hspace{1em}\cellcolor{gray!10}{Other Diagnosis} & \cellcolor{gray!10}{V1191} & \cellcolor{gray!10}{Binary (0/1)} & \cellcolor{gray!10}{Other lifetime mental health diagnosis} & \cellcolor{gray!10}{0 = No, 1 = Yes (2 → 0 recoded)} & \cellcolor{gray!10}{Psychiatric screening}\\
\hspace{1em}Alcoholism & V1265 & Binary (0/1) & Lifetime alcohol dependence diagnosis & 0 = No, 1 = Yes (2 → 0 recoded) & Substance use disorder assessment\\
\hspace{1em}\cellcolor{gray!10}{Armed Forces} & \cellcolor{gray!10}{V0023} & \cellcolor{gray!10}{Binary (0/1)} & \cellcolor{gray!10}{U.S. military service history} & \cellcolor{gray!10}{0 = No, 1 = Yes (2 → 0 recoded)} & \cellcolor{gray!10}{Military service verification}\\
\hspace{1em}Homeless & V0958 & Binary (0/1) & Housing instability prior to arrest & 0 = Stably housed, 1 = Homeless (2 → 0 recoded) & Pretrial housing assessment\\
\hspace{1em}\cellcolor{gray!10}{Children} & \cellcolor{gray!10}{V0982} & \cellcolor{gray!10}{Binary (0/1)} & \cellcolor{gray!10}{Has minor dependent children} & \cellcolor{gray!10}{0 = No, 1 = Yes (2 → 0 recoded)} & \cellcolor{gray!10}{Family responsibility assessment}\\
\hspace{1em}Mental Health Treatment & V1204 & Binary (0/1) & Received mental health treatment during incarceration & 0 = No treatment received post-admission, 1 = Received MHT during prison stay (2 → 0 recoded) & Prison mental health treatment records\\
\hspace{1em}\cellcolor{gray!10}{Alcohol/Drug Treatment} & \cellcolor{gray!10}{V1374} & \cellcolor{gray!10}{Binary (0/1)} & \cellcolor{gray!10}{Received substance abuse treatment during incarceration} & \cellcolor{gray!10}{0 = No treatment received post-admission, 1 = Received substance treatment during prison (2 → 0 recoded)} & \cellcolor{gray!10}{Prison substance abuse treatment records}\\*
\end{longtable}
\endgroup{}

\newpage

\section{Appendix B: Model Summary
Outputs}\label{appendix-b-model-summary-outputs}

\begin{verbatim}
## Stage 1 Logistic Model
\end{verbatim}

\begin{verbatim}
## 
## Call:
## glm(formula = at_risk ~ Years.of.Education + ADD.ADHD + Gender + 
##     offense_type + Children + Times.Arrested + Mental.Health.Treatment + 
##     Alcoholism, family = binomial(link = "logit"), data = cleaned2, 
##     na.action = na.exclude)
## 
## Coefficients:
##                           Estimate Std. Error z value Pr(>|z|)    
## (Intercept)              -0.105439   0.111897  -0.942 0.346044    
## Years.of.Education       -0.006716   0.007237  -0.928 0.353398    
## ADD.ADHD                  1.338996   0.047447  28.221  < 2e-16 ***
## GenderMale               -0.938654   0.047758 -19.654  < 2e-16 ***
## offense_typeOther         0.260657   0.350827   0.743 0.457494    
## offense_typeProperty      0.235999   0.067435   3.500 0.000466 ***
## offense_typePublic Order  0.193383   0.074804   2.585 0.009732 ** 
## offense_typeUnknown       0.217762   0.090668   2.402 0.016317 *  
## offense_typeViolent       0.120553   0.056052   2.151 0.031498 *  
## Children                 -0.086372   0.042271  -2.043 0.041025 *  
## Times.Arrested            0.013053   0.001667   7.829 4.94e-15 ***
## Mental.Health.Treatment   2.201666   0.050328  43.746  < 2e-16 ***
## Alcoholism                0.087787   0.038603   2.274 0.022961 *  
## ---
## Signif. codes:  0 '***' 0.001 '**' 0.01 '*' 0.05 '.' 0.1 ' ' 1
## 
## (Dispersion parameter for binomial family taken to be 1)
## 
##     Null deviance: 21322  on 15449  degrees of freedom
## Residual deviance: 16637  on 15437  degrees of freedom
##   (1708 observations deleted due to missingness)
## AIC: 16663
## 
## Number of Fisher Scoring iterations: 4
\end{verbatim}

\begin{verbatim}
## Stage 2 Logistic Model
\end{verbatim}

\begin{verbatim}
## 
## Call:
## glm.nb(formula = Number.Infractions ~ Race + num_diagnoses + 
##     Armed.Forces + +Times.Arrested + Years.of.Education + ADD.ADHD + 
##     Homeless + Children + Mental.Health.Treatment + Mental.Health.Treatment:num_diagnoses + 
##     Alcoholism + Alcohol.Drug.Treatment + Alcoholism:Alcohol.Drug.Treatment + 
##     Gender + offense_type, data = high_risk_group, init.theta = 0.225068483, 
##     link = log)
## 
## Coefficients:
##                                        Estimate Std. Error z value Pr(>|z|)    
## (Intercept)                           -0.656952   0.519424  -1.265 0.205953    
## RaceBlack                              1.035531   0.490403   2.112 0.034721 *  
## RaceHispanic                           1.054957   0.511488   2.063 0.039158 *  
## RaceMultiracial                        0.730125   0.489379   1.492 0.135714    
## RaceOther                              1.330617   0.536754   2.479 0.013175 *  
## RaceWhite                              0.439531   0.489300   0.898 0.369034    
## num_diagnoses                          0.072018   0.026254   2.743 0.006086 ** 
## Armed.Forces                          -0.145850   0.119248  -1.223 0.221300    
## Times.Arrested                         0.007912   0.001149   6.885 5.76e-12 ***
## Years.of.Education                    -0.037718   0.011595  -3.253 0.001142 ** 
## ADD.ADHD                               0.604736   0.063758   9.485  < 2e-16 ***
## Homeless                              -0.114640   0.113859  -1.007 0.314003    
## Children                              -0.395354   0.058881  -6.714 1.89e-11 ***
## Mental.Health.Treatment                0.324830   0.093639   3.469 0.000522 ***
## Alcoholism                             0.294180   0.081531   3.608 0.000308 ***
## Alcohol.Drug.Treatment                 0.023718   0.074325   0.319 0.749638    
## GenderMale                            -0.045294   0.062580  -0.724 0.469206    
## offense_typeOther                      0.400541   0.417414   0.960 0.337268    
## offense_typeProperty                  -0.014625   0.096488  -0.152 0.879527    
## offense_typePublic Order              -0.240554   0.111770  -2.152 0.031380 *  
## offense_typeUnknown                   -0.512680   0.133487  -3.841 0.000123 ***
## offense_typeViolent                    0.125330   0.085486   1.466 0.142624    
## num_diagnoses:Mental.Health.Treatment -0.014916   0.033334  -0.447 0.654547    
## Alcoholism:Alcohol.Drug.Treatment     -0.281589   0.109103  -2.581 0.009853 ** 
## ---
## Signif. codes:  0 '***' 0.001 '**' 0.01 '*' 0.05 '.' 0.1 ' ' 1
## 
## (Dispersion parameter for Negative Binomial(0.2251) family taken to be 1)
## 
##     Null deviance: 5716.9  on 7618  degrees of freedom
## Residual deviance: 5331.2  on 7595  degrees of freedom
##   (105 observations deleted due to missingness)
## AIC: 19762
## 
## Number of Fisher Scoring iterations: 1
## 
## 
##               Theta:  0.22507 
##           Std. Err.:  0.00655 
## 
##  2 x log-likelihood:  -19711.59400
\end{verbatim}

\newpage

\section{References}\label{references}

Prison Policy Initiative. (2026, April). Mental health: Research on the
prevalence and treatment of mental illness in the criminal legal system.
\url{https://www.prisonpolicy.org/research/mental_health/}

Anthropic. (2026). Claude (claude-sonnet-4-6) {[}Large language
model{]}. Data visualizations generated April 2026.
\url{https://claude.ai}

\end{document}
